\section{Slicing interpretation}
\label{sec:slicing-interpretation}

Earlier work attached approximation information at the syntactic level via a CBN translation into
the source language extended with a monad $\mathsf{Mon}$, modelling the per-constructor ``hole''
of prior Galois slicing work. That route does not extend to inductive types: the
substitution $\subst{\tau}{\sigma}{\alpha}$ in the typing rule for $\fold$ does not commute with a
syntactic Mon-tagging at constructor roots (the outer Mon lands at the wrong level once the
recursive variable is substituted). The obstruction is specific to source-to-source translations;
moving the Mon-tagging into the \emph{interpretation} avoids it, because semantic substitution is
just composition of interpretations and so commutes automatically.

This section gives such an interpretation, parameterised by a category $\cat{D}$ equipped with an
endofunctor $T : \cat{D} \to \cat{D}$ playing the role of $\mathsf{Mon}$. We leave the concrete
choice of $\cat{D}$ and $T$ for Galois slicing to a later section.

\subsection{Semantic assumptions}
\label{sec:slicing-interpretation:assumptions}

Assume a category $\cat{D}$ with:
\begin{itemize}
\item finite products $(\times, 1)$ and exponentials (so $\cat{D}$ is cartesian closed);
\item finite coproducts $(+, 0)$, with the products distributing over coproducts;
\item a strong monad $T : \cat{D} \to \cat{D}$ with unit $\eta : \id \Rightarrow T$, multiplication
  $\mu : T^2 \Rightarrow T$, and strength $\mathsf{st}_{X,Y} : X \times TY \to T(X \times Y)$;
\item a signature interpretation, in particular assigning each $\rho \in \PrimTy$ an object
  $\sem{\rho}_{\PrimTy} \in \cat{D}$;
\item initial algebras of the endofunctors arising from the interpretation of $\mu$-types
  (\secref{slicing-interpretation:mu}).
\end{itemize}

Commutativity of $T$ is not required: introductions package already-$T$-wrapped components inside
a fresh outer $T$ without sequencing the inner effects, and eliminations only ever fire one
monadic action at a time.

\subsection{Type interpretation}
\label{sec:slicing-interpretation:types}

A type $\tau$ kinded in context $\Delta = \alpha_1, \ldots, \alpha_n$ is interpreted as a functor
$\sem{\tau} : \cat{D}^n \to \cat{D}$. Closed types ($\Delta = \cdot$) thus interpret as objects of
$\cat{D}$. The interpretation puts $T$ at the root of each type constructor:
\begin{align*}
\sem{\alpha_i}                & = \pi_i \\
\sem{\rho}                    & = T\,\sem{\rho}_{\PrimTy} & \text{for $\rho \in \PrimTy$} \\
\sem{\tyUnit}                 & = T\,1 \\
\sem{\sigma \tyProd \tau}     & = T\,(\sem{\sigma} \times \sem{\tau}) \\
\sem{\sigma \tySum \tau}      & = T\,(\sem{\sigma} + \sem{\tau}) \\
\sem{\sigma \tyFun \tau}      & = T\,(\sem{\sigma} \Rightarrow \sem{\tau}) \\
\sem{\mu\alpha.\tau}          & = \mu(X \mapsto \sem{\tau}[X/\alpha]) & \text{see \secref{slicing-interpretation:mu}}
\end{align*}
where $\sem{\rho}_{\PrimTy}$ is the chosen object for $\rho$. The outer $T$ on primitives is not optional:
without it, eliminations whose result is a base type (e.g.\ $\tmFst M$ with $M : \rho \tyProd
\tau$) cannot be interpreted, since the Kleisli extension always produces a $T$-wrapped result
and the strong monad assumption alone gives no way to project the outer $T$ off.

The placement of $T$ at \emph{each} constructor root is what realises the per-constructor
approximation point that prior Galois slicing work attaches at every constructor. Substitution
commutes with the
interpretation:
\[
\sem{\subst{\tau}{\sigma}{\alpha}} = \sem{\tau}[\sem{\sigma}/\alpha]
\]
which is just function composition and so holds by definition.

\subsection{Typing contexts}
\label{sec:slicing-interpretation:contexts}

Contexts $\Gamma = x_1 : \tau_1, \ldots, x_n : \tau_n$ (with each $\tau_i$ closed) interpret as
\[
\sem{\Gamma} = \sem{\tau_1} \times \cdots \times \sem{\tau_n}
\]
using the cartesian product in $\cat{D}$. We do \emph{not} add a further $T$-wrap to context
entries: the type interpretation already places $T$ at the root of each $\tau_i$.

\subsection{Term interpretation}
\label{sec:slicing-interpretation:terms}

A term judgement $\Gamma \vdash t : \tau$ is interpreted as a morphism $\sem{t} : \sem{\Gamma} \to
\sem{\tau}$. Write $f^* = \mu \comp T(f) : TX \to TY$ for the Kleisli extension of $f : X \to TY$.
Introductions compose with $\eta$; eliminations take the Kleisli extension of the underlying
operation, using strength to thread $\Gamma$ where needed.

\begin{align*}
\sem{x_i}                              & = \pi_i \\
\sem{\tmUnit}                          & = \eta_{1} \comp {!_{\sem{\Gamma}}} \\
\sem{\tmPair{s}{t}}                    & = \eta_{\sem{\sigma}\times\sem{\tau}} \comp
                                            \prodM{\sem{s}}{\sem{t}} \\
\sem{\tmFst{s}}                        & = \pi_1^* \comp \sem{s} \\
\sem{\tmSnd{s}}                        & = \pi_2^* \comp \sem{s} \\
\sem{\tmInl{s}}                        & = \eta_{\sem{\sigma}+\sem{\tau}} \comp \mathsf{inj}_1
                                            \comp \sem{s} \\
\sem{\tmInr{s}}                        & = \eta_{\sem{\sigma}+\sem{\tau}} \comp \mathsf{inj}_2
                                            \comp \sem{s} \\
\sem{\tmCase{s}{x_1}{t_1}{x_2}{t_2}}   & = c^* \comp
                                            \mathsf{st}_{\sem{\Gamma}, \sem{\sigma_1}+\sem{\sigma_2}}
                                            \comp \prodM{\id}{\sem{s}} \\
\sem{\tmFun{x}{t}}                     & = \eta_{\sem{\sigma}\Rightarrow\sem{\tau}} \comp \lambda(\sem{t}) \\
\sem{\tmApp{s}{t}}                     & = \eval^* \comp
                                            \mathsf{st}_{\sem{\sigma}\Rightarrow\sem{\tau}, \sem{\sigma}}
                                            \comp \prodM{\sem{s}}{\sem{t}} \\
\sem{\phi(t_1, \ldots, t_n)}           & = (\eta_{\sem{\rho}_{\PrimTy}} \comp \sem{\phi}_{\Op})_{\diamond}
                                            \comp \prodM{\sem{t_1}}{\ldots, \sem{t_n}} \\
\sem{\tmRoll{t}}                       & = \mathsf{in}_F \comp \sem{t} \\
\sem{\syn{fold}\,s\,\syn{with}\,x.t}   & = h \comp \prodM{\id}{\sem{s}} & \text{see \secref{slicing-interpretation:mu}}
\end{align*}

In the $\syn{case}$ clause, $c : \sem{\Gamma} \times (\sem{\sigma_1}+\sem{\sigma_2}) \to \sem{\tau}$ is
the (Kleisli) arrow $\coprodM{\sem{t_1}}{\sem{t_2}} \comp \delta$ for $\delta$ the distributivity
isomorphism. The $\phi$ clause's iterated extension $(-)_{\diamond}$ is the obvious chained
Kleisli-extension that consumes each argument's outer $T$ in turn, threading $\Gamma$ via
strength.

\subsection{Inductive types}
\label{sec:slicing-interpretation:mu}

For $\Delta, \alpha \vdash \tau : \mathsf{type}$, the interpretation $\sem{\tau}$ is a functor
$\cat{D}^{|\Delta|+1} \to \cat{D}$, and so for any choice of objects for the variables in $\Delta$
gives an endofunctor $F = (X \mapsto \sem{\tau}[X/\alpha])$ on $\cat{D}$. We require initial
algebras for all such $F$. Then:
\[
\sem{\mu\alpha.\tau} = \mu F
\]
with algebra map $\mathit{in}_F : F(\mu F) \to \mu F$.

Substitution commutes by composition:
$\sem{\subst{\tau}{\sigma}{\alpha}} = F(\sem{\sigma})$, so the typing of $\tmRoll$ as
$\sem{\subst{\tau}{\mu\alpha.\tau}{\alpha}} \to \sem{\mu\alpha.\tau}$ becomes
$F(\mu F) \to \mu F$, which is exactly $\mathit{in}_F$.

For $\fold\,s\,\with\,x.t$ with $s : \mu\alpha.\tau$ and $x : \subst{\tau}{\sigma}{\alpha} \vdash t
: \sigma$ (writing $\sigma$ for the result type), the body $t$ interprets as
\[
\sem{t} : \sem{\Gamma} \times F(\sem{\sigma}) \to \sem{\sigma}
\]
which curries to a map $\sem{\Gamma} \to F(\sem{\sigma}) \Rightarrow \sem{\sigma}$, i.e.\ a
parameterised $F$-algebra. The unique algebra morphism $h : \mu F \to \sem{\sigma}$ exists by
initiality (in fact, a parameterised version of it, using the cartesian closed structure to handle
the $\Gamma$-context).

Crucially, no monad-substitution commutation lemma is required: the equation
$\sem{\subst{\tau}{\sigma}{\alpha}} = F(\sem{\sigma})$ that the typing of $\fold$ demands is
satisfied definitionally by the type interpretation.

\subsection{Where the approximation lives}
\label{sec:slicing-interpretation:approximation}

The monad $T$ is the carrier of the per-constructor approximation. Its concrete realisation
depends on which analysis we are doing:
\begin{itemize}
\item For Galois slicing, $T$ adds a ``hole'' value or top lattice element at each constructor.
\item For automatic differentiation, $T$ would attach tangent-space information.
\end{itemize}
A subsequent section will spell out $T$ concretely for Galois slicing in $\cat{D} = \Fam(\cat{C})$
for an appropriate $\cat{C}$, and verify the strong monad laws there. The point of the present
section is that, with such a $T$ assumed, the interpretation is uniform across all type
constructors including $\mu$, and no substitution-commutation obligation arises at the syntactic
level.
